\section{Analisi degli aspetti di sicurezza}\label{sec:sicurezza}

L'analisi di sicurezza copre i principali pilastri del modello CIA
(\emph{Confidentiality, Integrity, Availability}) ed è suddivisa fra
\emph{controlli già implementati nel codice} e \emph{controlli previsti
come future work} (rinviati alla fase di produzione).

\subsection{Modello di minaccia}
\begin{table}[H]
    \centering
    \begin{tabularx}{\textwidth}{@{}p{4.5cm}X@{}}
        \toprule
        \textbf{Minaccia} & \textbf{Mitigazione adottata} \\
        \midrule
        Credential stuffing &
        Hashing password lento (PBKDF2-SHA256, $\sim 600k$ iterazioni di
        default) $+$ rate limit (deferred). \\
        User enumeration &
        Risposte unificate su login fallito (\code{InvalidCredentialsError},
        HTTP 401) anche quando l'utente non esiste. \\
        Timing attack &
        \code{check\_password\_hash} di werkzeug è constant-time e viene
        chiamato anche su utenti inesistenti. \\
        Replay del token JWT &
        Scadenza breve dell'\emph{access token} (30 min); \emph{refresh
        token} revocabile lato server come future work (blacklist Redis). \\
        Spam di pubblicazioni anonime &
        Sfida Proof of Work obbligatoria; token consumato in maniera
        atomica via \code{UPDATE WHERE used = 0} con check del
        \code{rowcount}. \\
        Privilege escalation &
        Le route \endpoint{PUT}{/api/songs/<id>} ed
        \endpoint{DELETE}{/api/songs/<id>} verificano che
        \code{song.user\_id == current\_user\_id}. \\
        Attacchi DoS via payload abnorme &
        \code{MAX\_CONTENT\_LENGTH = 256\,KB} a livello WSGI; limiti
        per-campo (50~KB \code{plainLyrics}, 100~KB \code{syncedLyrics},
        300 char titolo) a livello applicativo. \\
        SQL injection &
        Uso esclusivo dell'ORM SQLAlchemy con parametri \emph{bound}; nessun
        \code{text()} o concatenazione di stringhe. \\
        XSS / CSRF &
        API stateless con token in header \code{Authorization} (non
        cookie); CSRF non applicabile. XSS prevenuto da rendering React
        Native (no \code{innerHTML}). \\
        \bottomrule
    \end{tabularx}
    \caption{Modello di minaccia e mitigazioni.}
\end{table}

\subsection{Autenticazione e gestione utenti}
\begin{itemize}
    \item Password con hashing \textbf{PBKDF2-SHA256} e salt randomico di
    16 byte (defaults di \code{werkzeug.security.generate\_password\_hash}).
    \item Validazione lato server: username
    \verb|^[A-Za-z0-9_]{3,30}$|, email regex, password $\geq 8$ caratteri
    con almeno una cifra.
    \item Univocità di username ed email è verificata case-insensitive.
    \item Il login accetta indistintamente username o email come
    identificatore.
\end{itemize}

\subsection{Crittografia delle comunicazioni}
\begin{itemize}
    \item In \textbf{sviluppo}: HTTP su rete locale.
    \item In \textbf{produzione} (deferred): TLS terminato da reverse proxy
    (\emph{nginx} o \emph{Caddy}) con certificato \emph{Let's Encrypt} e
    header \code{Strict-Transport-Security: max-age=31536000; includeSubDomains}.
\end{itemize}

\subsection{Protezione dei dati a riposo}
\begin{itemize}
    \item Le password sono memorizzate solo come hash; il valore in chiaro
    non viene mai loggato né serializzato.
    \item Il database SQLite vive in un singolo file (\code{lyrics.db}); in
    produzione si raccomanda cifratura del filesystem o uso di SQLCipher.
    \item I JWT non vengono persistiti lato server: sono \emph{stateless} e
    autocontenuti.
\end{itemize}

\subsection{Gestione delle vulnerabilità}
\begin{itemize}
    \item Dipendenze pinnate in \code{requirements.txt} (Python) e
    \code{package.json} (Node).
    \item Audit periodico previsto con \code{pip-audit} e \code{npm audit}
    (deferred).
\end{itemize}

\subsection{Backup e ripristino}
\begin{itemize}
    \item \textbf{Sviluppo:} nessun backup automatico; il file SQLite è
    contenuto nel volume del progetto.
    \item \textbf{Produzione (deferred):} snapshot giornaliero del file
    \code{lyrics.db} con retention di 7 giorni; copia offsite via
    \code{rsync} verso storage cifrato.
\end{itemize}

\subsection{Monitoraggio e logging}
\begin{itemize}
    \item \textbf{Sviluppo:} log su stdout dal Werkzeug dev server.
    \item \textbf{Produzione (deferred):} logging strutturato in JSON con
    rotazione (\emph{RotatingFileHandler}) e aggregatore esterno
    (Loki/ELK). Eventi tracciati: login falliti, pubblicazioni, mutazioni
    su risorse.
\end{itemize}

\subsection{Conformità legale}
\subsubsection{GDPR (Regolamento UE 2016/679)}
\begin{itemize}
    \item \textbf{Dati personali raccolti:} username, email, hash della
    password, data di creazione.
    \item \textbf{Base giuridica:} consenso esplicito al momento della
    registrazione.
    \item \textbf{Diritto di accesso:} fornito tramite l'endpoint
    \endpoint{GET}{/auth/me}.
    \item \textbf{Diritto all'oblio:} previsto come estensione futura
    (\endpoint{DELETE}{/auth/me}) con cancellazione del record User e
    anonimizzazione dei brani associati (\code{user\_id = NULL}).
    \item \textbf{Cifratura}: in produzione tutti i dati personali
    transitano cifrati via TLS.
\end{itemize}

\subsubsection{Diritto d'autore sui testi}
I testi pubblicati possono essere soggetti a copyright. Il sistema si
qualifica come \emph{piattaforma di hosting di contenuti utente} e adotta
le seguenti misure (alcune deferred):
\begin{itemize}
    \item Termini di servizio espliciti che vietano l'upload di materiale
    protetto di cui l'utente non detenga i diritti.
    \item Endpoint di \emph{takedown} per la rimozione su segnalazione del
    titolare dei diritti (deferred).
    \item Attribuzione dei contributi tramite il campo \code{submittedBy}
    per rendere tracciabili gli autori dei caricamenti.
\end{itemize}

\subsection{Mitigazioni implementate (riepilogo)}
\begin{enumerate}
    \item Hash password PBKDF2-SHA256 con salt.
    \item Risposte uniformi sul login fallito (no user enumeration).
    \item Controllo della password con tempo costante anche per utenti
    inesistenti.
    \item JWT firmati HS256 con scadenza differenziata access/refresh.
    \item Proof of Work one-shot atomico (\code{UPDATE WHERE used = 0}).
    \item Limite globale del body (\code{MAX\_CONTENT\_LENGTH = 256 KB}).
    \item Limiti dimensionali per campo (50~KB plain, 100~KB synced,
    300~char titolo, 0--7200~s durata).
    \item Autorizzazione su \code{PUT}/\code{DELETE} dei brani con check di
    proprietà.
    \item ORM SQLAlchemy con parametri bound (no SQL injection).
    \item Validazione regex su username/email/password lato server.
    \item Error handler globali 404/405/500 senza leak di stack trace.
\end{enumerate}

\subsection{Mitigazioni deferred to production}
\begin{enumerate}
    \item TLS terminator e HSTS.
    \item Rate limiting su login/register/publish/request-challenge.
    \item Token blacklist server-side per logout effettivo.
    \item Cookie \code{httpOnly} + \code{Secure} + \code{SameSite=Strict}
    per il refresh token (in alternativa allo storage client).
    \item Logging strutturato e aggregatore esterno.
    \item Backup periodico cifrato del DB.
    \item Cifratura del filesystem (o SQLCipher).
    \item Endpoint \endpoint{DELETE}{/auth/me} per il diritto all'oblio
    GDPR.
    \item Endpoint \emph{takedown} per la rimozione di contenuti su
    segnalazione.
    \item Audit periodico delle dipendenze con \code{pip-audit} e
    \code{npm audit}.
\end{enumerate}
